\chapter{Consensus and Formation Control}
\label{ch:ch23}
% Function names for pseudocode are prefixed with "Consensus" to stay unique
% across the book.
\SetKwFunction{ConsensusFormationStep}{FormationStep}
\SetKwFunction{ConsensusSecondOrderStep}{SecondOrderFormationStep}

\begin{objectives}
\begin{itemize}
  \item Build the communication graph of a swarm from the drone positions and the radio range, write down its adjacency, degree and Laplacian matrices, and read connectivity and convergence speed off the Laplacian spectrum.
  \item State the consensus protocol $\dot{\vect{x}} = -\mat{L}\vect{x}$, prove that on a connected graph it drives every drone to the average of the initial states at the rate $\lambda_2$, and choose a safe step size for its discrete-time version.
  \item Turn consensus into a formation controller by prescribing displacements, pin a leader to a reference path, and explain why the formation is reached only up to a translation.
  \item Define the formation error and the algebraic connectivity as the two quantities that a formation-flying swarm monitors, and express the communication radius as a hard constraint or as a potential.
  \item Extend the controller to double-integrator drones and to flocking with separation, alignment and navigation terms, and implement and test all of it in \python.
\end{itemize}
\end{objectives}

% ---------------------------------------------------------------------
\section{Why this matters for a drone swarm}
\label{sec:ch23-motivation}

Four camera drones film a moving vehicle from four sides. The images are only
useful if the drones keep a square of side $2$~m around the target, and the
drones can only coordinate if every one of them stays within the radio range
$R_{\mathrm{comm}}$ of at least one other drone. The global planner of the
hybrid architecture (\cref{ch:ch24}) has produced a conflict-free path for the
square as a whole with \cbs or \ecbs (\cref{ch:ch09,ch:ch10}); the reactive
layer (\orca, \cref{ch:ch13}; DWA, \cref{ch:ch14}) will push single drones off
that path when a bird crosses the scene. Between these two layers something has
to hold the square together: a controller that lets each drone use only what
it can hear from its neighbours, that keeps the relative positions, that keeps
the communication graph connected, that yields to an avoidance manoeuvre and
that pulls the square back together afterwards. That controller is the subject
of this chapter, and its two constraints, formation preservation and the
communication radius, are the last two constraints of the capstone system of
\cref{ch:ch24}; \cref{ch:ch25} measures how well they are met with the
formation error and the communication violations defined here.

The controller is built from a single rule, \textbf{consensus}: every drone
moves its state towards the average of its neighbours' states. Consensus is
one of the few results in this book with a complete, short and beautiful
theory. Its convergence is decided by one number, the second-smallest
eigenvalue $\lambda_2$ of the graph Laplacian, which is positive exactly when
the communication graph is connected and which is also the rate at which the
drones agree. Everything else in the chapter, leader--follower tracking,
displacement-based formation control, second-order control for
double-integrator drones and flocking, is consensus applied to a shifted or
extended state.

The chapter follows the usual order: the rule and its picture, the graph
matrices and the two error metrics, the controllers (consensus, its
discrete-time form, a pinned leader, formations, the second-order version), a
worked example of four drones, the convergence theorems, and finally flocking,
connectivity maintenance and the implementation. Everything after
\cref{sec:ch23-algorithm} rests on the single observation of
\cref{def:ch23-formation}: a formation is consensus on a shifted state.

% ---------------------------------------------------------------------
\section{The idea in plain words}
\label{sec:ch23-intuition}

Imagine that each drone holds a number, say its altitude, and that the swarm
wants to agree on a common altitude without a leader and without anybody
knowing all the numbers. The rule is: at every instant, move your number
towards the numbers of the drones you can hear (\cref{fig:ch23-idea}). A drone
whose neighbours are all higher climbs, a drone in the middle of its
neighbours stays. Two things follow immediately. Whatever one drone gains, its
neighbour loses, so the \emph{sum} of the numbers never changes; and the
\emph{spread} of the numbers can only shrink, because the largest number can
only go down and the smallest can only go up. If the communication graph is
connected, the spread shrinks to zero and the swarm settles on the average of
the initial numbers. How fast it does so depends on how well the graph is
connected: a long chain of drones agrees slowly, a dense cluster agrees fast.

\begin{figure}[tb]
  \centering
  \inputfigure{ch23/idea}
  \caption{Left: the communication graph of six drones with radio range
  $R_{\mathrm{comm}}$; drone~3 hears drones 2, 4 and 5 and nobody else, so its
  update uses three neighbours. Right: the consensus rule for one drone. Each
  grey arrow is the difference $\vect{x}_j - \vect{x}_i$ to a neighbour; their
  sum (orange) is $\abs{N_i}$ times the vector from $\vect{x}_i$ to the
  neighbours' average, so drone $i$ moves towards the average of the drones it
  can hear.}
  \label{fig:ch23-idea}
\end{figure}

Formation control is the same rule applied to a shifted number. Give every
drone a slot $\vect{o}_i$ in a template of the formation, say the four corners
of a square, and let the drones run consensus not on their positions but on
$\pos_i - \vect{o}_i$, the position of the template's origin as drone $i$ sees
it. When they agree on that origin, every drone is in its slot and the swarm
is in formation. The template can be placed anywhere: consensus decides
\emph{where}, the offsets decide the \emph{shape}. If one drone is told where
the origin should be, for example by the \cbs path, the whole formation
follows it.

\begin{keyidea}
Consensus is the rule $\dot{\vect{x}}_i = \sum_{j \in N_i} (\vect{x}_j -
\vect{x}_i)$, in matrix form $\dot{\vect{x}} = -\mat{L}\vect{x}$ with the
graph Laplacian $\mat{L}$. On a connected graph it drives every drone to the
average of the initial states, and the disagreement decays like
$e^{-\lambda_2 t}$, where $\lambda_2$ is the algebraic connectivity of the
graph. A formation is consensus on the shifted positions $\pos_i - \vect{o}_i$:
the drones agree on where the template sits, and the offsets give the shape.
\end{keyidea}

% ---------------------------------------------------------------------
\section{Problem statement and notation}
\label{sec:ch23-problem}

\subsection{The communication graph and its matrices}

\begin{definition}[Communication graph]
\label{def:ch23-graph}
A swarm of $n$ drones with positions $\pos_1, \dots, \pos_n \in \R^m$
($m = 2$ or $3$) forms the \textbf{communication graph}\index{communication graph}
$G = (V, E)$ with $V = \set{1, \dots, n}$ and an edge $\set{i, j} \in E$ whenever
drones $i$ and $j$ can exchange messages. For a radio of range
$R_{\mathrm{comm}}$ this is the \textbf{radius graph}\index{radius graph}
\begin{equation}
  \set{i, j} \in E \iff \norm{\pos_i - \pos_j} \le R_{\mathrm{comm}}, \qquad i \ne j.
  \label{eq:ch23-radius}
\end{equation}
The \textbf{neighbourhood}\index{neighbourhood $N_i$} of drone $i$ is
$N_i = \set{j \given \set{i, j} \in E}$ and its \textbf{degree} is
$d_i = \abs{N_i}$. The graph is \textbf{directed} when the links are one-way:
then $(j, i) \in E$ means that $i$ receives from $j$, and $N_i$ is the set of
drones that $i$ hears. A graph that changes with time is a \textbf{switching
topology}\index{switching topology}; the radius graph switches whenever a
pair of drones crosses the distance $R_{\mathrm{comm}}$.
\end{definition}

The mathematics refresher (\cref{ch:appB}) collects the linear algebra used
below: symmetric matrices have real eigenvalues and an orthonormal basis of
eigenvectors, and a symmetric matrix is positive semidefinite exactly when
its quadratic form is non-negative.

\begin{definition}[Adjacency, degree and Laplacian matrices]
\label{def:ch23-laplacian}
The \textbf{adjacency matrix}\index{adjacency matrix} $\mat{A} \in \R^{n
\times n}$ has $a_{ij} = 1$ if $\set{i, j} \in E$ and $a_{ij} = 0$ otherwise
(in particular $a_{ii} = 0$); weighted graphs use $a_{ij} > 0$ for the strength
of the link. The \textbf{degree matrix}\index{degree matrix} is
$\mat{D} = \diag(d_1, \dots, d_n)$ with $d_i = \sum_j a_{ij}$, and the
\textbf{graph Laplacian}\index{graph Laplacian}\index{Laplacian|see{graph Laplacian}} is
\begin{equation}
  \mat{L} = \mat{D} - \mat{A}, \qquad
  \ell_{ij} = \begin{cases} d_i & i = j,\\ -a_{ij} & i \ne j. \end{cases}
  \label{eq:ch23-laplacian}
\end{equation}
Its eigenvalues are written in increasing order,
$\lambda_1 \le \lambda_2 \le \dots \le \lambda_n$; $\lambda_2$ is the
\textbf{algebraic connectivity}\index{algebraic connectivity $\lambda_2$}
of the graph \cite{fiedler1973algebraic}. $\vect{1}$ denotes the vector of $n$
ones.
\end{definition}

\begin{proposition}[Properties of the Laplacian of an undirected graph]
\label{prop:ch23-laplacian}
Let $G$ be undirected with $n$ vertices and let $d_{\max}$ and $d_{\min}$ be its
largest and smallest degree. Then
\begin{enumerate}[label=(\roman*)]
  \item every row of $\mat{L}$ sums to zero: $\mat{L}\vect{1} = \vect{0}$ and
    $\vect{1}\T \mat{L} = \vect{0}\T$;
  \item $\mat{L}$ is symmetric and
    $\vect{x}\T \mat{L} \vect{x} = \sum_{\set{i,j} \in E} a_{ij} (x_i - x_j)^2 \ge 0$
    for every $\vect{x} \in \R^n$, so $\mat{L}$ is positive semidefinite and
    $0 = \lambda_1 \le \lambda_2 \le \dots \le \lambda_n$;
  \item the multiplicity of the eigenvalue $0$ equals the number of connected
    components of $G$; in particular $\lambda_2 > 0$ if and only if $G$ is
    connected;
  \item $\lambda_n \le 2 d_{\max}$;
  \item if $G$ is not complete then $\lambda_2 \le d_{\min}$, and the complete
    graph $K_n$ has $\lambda_2 = n$.
\end{enumerate}
\end{proposition}
\begin{proof}
(i) The diagonal entry $d_i$ of row $i$ equals the sum of the $a_{ij}$ that are
subtracted in the same row. (ii) Symmetry is inherited from $\mat{A}$. Expand
$\vect{x}\T \mat{L} \vect{x} = \sum_i d_i x_i^2 - \sum_{i \ne j} a_{ij} x_i x_j$
and group the terms of each edge: $a_{ij}(x_i^2 + x_j^2 - 2 x_i x_j) =
a_{ij}(x_i - x_j)^2$. A symmetric matrix with a non-negative quadratic form
has non-negative eigenvalues, and $\lambda_1 = 0$ by (i). (iii) By (ii),
$\mat{L}\vect{x} = \vect{0}$ holds if and only if $\vect{x}\T \mat{L} \vect{x} =
0$, that is, $x_i = x_j$ on every edge, that is, $\vect{x}$ is constant on
every connected component. The indicator vectors of the components therefore
span the null space, whose dimension is the number of components. (iv) By
Gershgorin's theorem every eigenvalue lies in a disc centred at some diagonal
entry $d_i$ with radius $\sum_{j \ne i} \abs{\ell_{ij}} = d_i$, hence in
$[0, 2 d_{\max}]$. (v) is Fiedler's theorem \cite{fiedler1973algebraic}: for a
non-complete graph $\lambda_2$ is at most the vertex connectivity, which is at
most $d_{\min}$. For $K_n$, $\mat{L} = n\mat{I} - \vect{1}\vect{1}\T$ has the
eigenvalue $n$ on the whole space orthogonal to $\vect{1}$.
\end{proof}

Two examples fix the scale of $\lambda_2$. A path of $n$ drones in a line, each
hearing only its two neighbours, has $\lambda_2 = 2(1 - \cos(\pi/n))$, which is
$0.27$ for $n = 6$ and shrinks like $\pi^2/n^2$; a complete graph on the same
six drones has $\lambda_2 = 6$. Connectivity is a matter of degree, and
$\lambda_2$ measures it on a continuous scale from $0$ (disconnected) upwards.

\subsection{Consensus}

\begin{definition}[Consensus]
\label{def:ch23-consensus}
Each drone $i$ holds a state $\vect{x}_i \in \R^m$ (a scalar, a position, a
velocity, an estimate). The swarm \textbf{reaches consensus}\index{consensus!definition}
if $\vect{x}_i(t) - \vect{x}_j(t) \to \vect{0}$ for all pairs $i, j$ as
$t \to \infty$, and \textbf{average consensus}\index{consensus!average} if the
common limit is the average $\bar{\vect{x}} = \frac{1}{n}\sum_i \vect{x}_i(0)$
of the initial states. Stacking the states of a scalar problem into
$\vect{x} = (x_1, \dots, x_n)\T$, the \textbf{disagreement vector}\index{disagreement vector}
is $\vect{\delta} = \vect{x} - \bar{x}\vect{1}$, the component of $\vect{x}$
orthogonal to $\vect{1}$; the swarm is in consensus exactly when
$\vect{\delta} = \vect{0}$. Vector states are handled one coordinate at a time.
\end{definition}

\subsection{Formations and the formation error}

\begin{definition}[Displacement-based formation]
\label{def:ch23-formation}
A \textbf{formation}\index{formation!template} is described by a template of
offsets $\vect{o}_1, \dots, \vect{o}_n \in \R^m$ in a common frame, one slot per
drone. The \textbf{desired displacement}\index{formation!desired displacement $\vect{d}_{ij}$}
from drone $i$ to drone $j$ is
\begin{equation}
  \vect{d}_{ij} = \vect{o}_j - \vect{o}_i, \qquad\text{so that}\qquad
  \vect{d}_{ji} = -\vect{d}_{ij}
  \quad\text{and}\quad
  \vect{d}_{ij} + \vect{d}_{jk} + \vect{d}_{ki} = \vect{0}.
  \label{eq:ch23-dij}
\end{equation}
The \textbf{formation graph}\index{formation!graph} $G_F = (V, E_F)$ lists the
pairs whose displacement is prescribed and measured; it is connected and is
usually the communication graph, a spanning subgraph of it, or the complete
graph. The swarm is \textbf{in formation} if $\pos_j - \pos_i = \vect{d}_{ij}$
for every $\set{i, j} \in E_F$; because $G_F$ is connected this holds if and
only if $\pos_i = \vect{o}_i + \vect{c}$ for all $i$ and one common translation
$\vect{c}$.
\end{definition}

A formation is therefore defined only up to translation, which is exactly what
the leader supplies. The shifted variable $\vect{y}_i = \pos_i - \vect{o}_i$
is the position of the template origin as seen by drone $i$, and the swarm is
in formation if and only if all $\vect{y}_i$ agree. This is the observation
that turns consensus into formation control in \cref{sec:ch23-formation}.

\begin{definition}[Formation error]
\label{def:ch23-error}
The \textbf{formation error}\index{formation error} of the positions
$\pos = (\pos_1, \dots, \pos_n)$ with respect to the formation
$(\vect{o}_i)$ and the formation graph $G_F$ is the root-mean-square violation
of the desired displacements over the edges of $G_F$,
\begin{equation}
  e_F(\pos) = \sqrt{\frac{1}{\abs{E_F}} \sum_{\set{i,j} \in E_F}
    \norm{\pos_j - \pos_i - \vect{d}_{ij}}^2}
  = \sqrt{\frac{1}{\abs{E_F}} \sum_{\set{i,j} \in E_F}
    \norm{\vect{y}_j - \vect{y}_i}^2}.
  \label{eq:ch23-error}
\end{equation}
The \textbf{centred formation error} removes the best translation first:
with $\bar{\vect{y}} = \frac{1}{n}\sum_i \vect{y}_i$,
\begin{equation}
  e_c(\pos) = \sqrt{\frac{1}{n} \sum_{i=1}^{n} \norm{\vect{y}_i - \bar{\vect{y}}}^2}.
  \label{eq:ch23-errorc}
\end{equation}
\end{definition}

Both errors have the unit of a length, both are invariant under a translation
of the whole swarm, and both are zero exactly when the swarm is in formation.
$e_F$ is the metric used by \cref{ch:ch25}: it only needs the pairs of the
formation graph, so a drone can compute its own contribution from its
neighbours' messages. The translation $\bar{\vect{y}}$ that $e_c$ removes is
the least-squares fit of the template to the current positions, and for the
complete formation graph the two metrics are proportional,
$e_F^2 = \frac{2n}{n-1}\, e_c^2$ (\cref{exr:ch23-average}). Two more
quantities are monitored alongside $e_F$: the \textbf{edge lengths}
$\ell_{ij} = \norm{\pos_i - \pos_j}$ of the edges that must stay in range, and
the number of \textbf{communication violations}\index{communication violation},
the edges of a chosen set $E_{\mathrm{keep}} \subseteq E_F$ that are longer than
$R_{\mathrm{comm}}$ at a given time.

% ---------------------------------------------------------------------
\section{The algorithm: consensus and its extensions}
\label{sec:ch23-algorithm}

\subsection{The consensus protocol}
\label{sec:ch23-protocol}

The \textbf{consensus protocol}\index{consensus!protocol} of
\textcite{olfatisaber2004consensus} is the rule of \cref{sec:ch23-intuition}
written as a differential equation. Drone $i$ moves its state with the velocity
\begin{equation}
  \dot{\vect{x}}_i = \sum_{j \in N_i} (\vect{x}_j - \vect{x}_i)
  = \sum_{j=1}^{n} a_{ij} (\vect{x}_j - \vect{x}_i), \qquad i = 1, \dots, n.
  \label{eq:ch23-protocol}
\end{equation}
The right-hand side is $-d_i \vect{x}_i + \sum_j a_{ij} \vect{x}_j$, which is
row $i$ of $-\mat{L}\vect{x}$, so the whole swarm obeys the linear system
\begin{equation}
  \dot{\vect{x}} = -\mat{L}\vect{x}.
  \label{eq:ch23-matrix}
\end{equation}
Three features of \cref{eq:ch23-protocol} matter for a swarm. The rule is
\emph{local}: drone $i$ needs only the states of the drones it hears. It is
\emph{leaderless} and \emph{symmetric}: no drone is special and every message
is used the same way. And it \emph{conserves the sum}: by
\cref{prop:ch23-laplacian}(i), $\frac{d}{dt}(\vect{1}\T\vect{x}) =
-\vect{1}\T\mat{L}\vect{x} = 0$, so the average $\bar{x}$ is a constant of the
motion. Every drone that runs \cref{eq:ch23-protocol} is therefore computing
the same number, the average of the initial states, without anybody ever
collecting them. A \textbf{consensus gain}\index{consensus!gain $k_c$} $k_c > 0$
in front of the sum only rescales time ($\vect{x}(t)$ becomes
$\vect{x}(k_c t)$) and is omitted in the analysis. The letter $k$ is reserved
for the time step and the mode index.

\subsection{Discrete-time consensus and the step size}
\label{sec:ch23-discrete}

A drone updates its command at a control period $\dt$, so what it really runs
is the Euler step of \cref{eq:ch23-matrix} with the step size
$\eps = k_c\,\dt$:
\begin{equation}
  \vect{x}_{k+1} = (\mat{I} - \eps\mat{L})\,\vect{x}_k, \qquad\text{that is}\qquad
  \vect{x}_{i,k+1} = \vect{x}_{i,k} + \eps \sum_{j \in N_i} (\vect{x}_{j,k} - \vect{x}_{i,k}).
  \label{eq:ch23-discrete}
\end{equation}
The matrix $\mat{P} = \mat{I} - \eps\mat{L}$ is called the Perron matrix of
the graph \cite{olfatisaber2007consensus}. Its rows sum to one, and its entries
$1 - \eps d_i$ on the diagonal and $\eps a_{ij}$ off it are all non-negative if
and only if $\eps \le 1/d_{\max}$. In that case every drone's new state is a
weighted average of its own and its neighbours' states, in which it keeps the
weight $1 - \eps d_i$ on itself, and the maximum can never grow nor the minimum
shrink. The \textbf{step-size condition}\index{consensus!step-size condition}
\begin{equation}
  0 < \eps < \frac{1}{d_{\max}}
  \label{eq:ch23-stepsize}
\end{equation}
is the sufficient condition every drone can check from local information (it
knows its own degree, and a bound on the largest degree is a design parameter
of the swarm). The exact condition is $\eps < 2/\lambda_n$, which by
\cref{prop:ch23-laplacian}(iv) is implied by \cref{eq:ch23-stepsize}; both are
proved in \cref{thm:ch23-discrete}. Above the exact bound the iteration
diverges, with the fastest mode of the graph flipping sign at every step; the
worked example shows this happening.

\subsection{Leader--follower consensus}
\label{sec:ch23-leader}

Average consensus agrees on a value that nobody chose. To make the swarm
follow a reference $\vect{r}(t)$, for example the \cbs path of the formation,
one or more drones are \textbf{pinned}\index{consensus!leader--follower}\index{pinning}
to it \cite{ren2008distributed}:
\begin{equation}
  \dot{\vect{x}}_i = \sum_{j \in N_i} (\vect{x}_j - \vect{x}_i)
  + b_i\, k_l\, (\vect{r} - \vect{x}_i) + \dot{\vect{r}}, \qquad
  b_i = \begin{cases} 1 & \text{drone $i$ is pinned,}\\ 0 & \text{otherwise,} \end{cases}
  \label{eq:ch23-pinned}
\end{equation}
with a pinning gain $k_l > 0$. The pinned drones are the
\textbf{leaders}\index{leader--follower}; the others follow them through the
graph and never see $\vect{r}$. With $\mat{B} = \diag(b_1, \dots, b_n)$ the
tracking error $\vect{e} = \vect{x} - \vect{1}\otimes\vect{r}$ (every drone
compared with the reference) obeys
\begin{equation}
  \dot{\vect{e}} = -(\mat{L} + k_l \mat{B})\,\vect{e},
  \label{eq:ch23-pinnedmatrix}
\end{equation}
and \cref{prop:ch23-leader} shows that $\mat{L} + k_l\mat{B}$ is positive
definite whenever the graph is connected and at least one drone is pinned, so
that $\vect{e} \to \vect{0}$. The \textbf{feed-forward}\index{feed-forward}
term $\dot{\vect{r}}$ in \cref{eq:ch23-pinned} is the velocity of the reference,
which the leader knows from its path and the followers can receive from it
through the graph. Without it the followers still converge, but to a
formation that trails the reference by a constant lag proportional to the speed
(\cref{exr:ch23-secondorder}); \cref{fig:ch23-leader} shows the difference.

\subsection{Formation control by displacements}
\label{sec:ch23-formation}

The \textbf{displacement-based formation controller}\index{formation control!displacement-based}
\cite{oh2015survey,ren2008distributed} replaces the difference of the states
by the difference minus the desired displacement:
\begin{equation}
  \dot{\pos}_i = \sum_{j \in N_i} (\pos_j - \pos_i - \vect{d}_{ij}).
  \label{eq:ch23-formation}
\end{equation}
Drone $i$ moves towards the point $\pos_j - \vect{d}_{ij}$ where it should be
according to each neighbour $j$, and it stops when it is there for all of them.
With $\vect{d}_{ij} = \vect{o}_j - \vect{o}_i$ and the shifted variables
$\vect{y}_i = \pos_i - \vect{o}_i$ the summand is $\vect{y}_j - \vect{y}_i$,
so that
\begin{equation}
  \dot{\vect{y}}_i = \sum_{j \in N_i} (\vect{y}_j - \vect{y}_i),
  \qquad \dot{\vect{y}} = -\mat{L}\vect{y} \text{ in every coordinate}:
  \label{eq:ch23-shift}
\end{equation}
the shifted variables run the plain consensus protocol
\cref{eq:ch23-matrix}. The offsets never appear in the dynamics of
$\vect{y}$; they only decide what the agreement means. Consequently every
theorem about consensus is a theorem about formation control: on a connected
graph the $\vect{y}_i$ converge to their initial average
$\bar{\vect{y}} = \frac{1}{n}\sum_i (\pos_i(0) - \vect{o}_i)$, and the drones
settle at $\pos_i = \vect{o}_i + \bar{\vect{y}}$, the template translated by the
initial centroid shift (\cref{thm:ch23-formation}). Pinning a leader to
$\vect{r} + \vect{o}_i$ as in \cref{eq:ch23-pinned} pins the template origin to
$\vect{r}$ and makes the formation follow the path.

\begin{pitfall}[Inconsistent displacement vectors]
Every $\vect{d}_{ij}$ must satisfy \cref{eq:ch23-dij}. If drone $i$ wants $j$
two metres to its east but $j$ wants $i$ one metre to its west, the pair
never settles: the sum of the displacement terms over the graph is no longer
zero, and the centroid of the swarm drifts with a constant velocity while the
relative positions converge to a compromise (\cref{exr:ch23-inconsistent}).
The same happens when the vectors do not close around a cycle,
$\vect{d}_{12} + \vect{d}_{23} + \vect{d}_{31} \ne \vect{0}$: no set of
positions satisfies all edges, the controller minimises the residual and a
formation error remains for ever. The safe way is to store the offsets
$\vect{o}_i$ and to compute $\vect{d}_{ij} = \vect{o}_j - \vect{o}_i$ on the
fly, which makes both conditions automatic; never store the $\vect{d}_{ij}$
themselves.
\end{pitfall}

\Cref{alg:ch23-formation} is one control period of drone $i$ with all of the
above combined. It is written from the point of view of a single drone, because
that is how it runs: there is no central computer that assembles
$\mat{L}$.

\begin{algorithm}[htb]
\caption{One control period of the displacement-based formation controller on drone $i$}
\label{alg:ch23-formation}
\KwIn{own position $\pos_i$ and offset $\vect{o}_i$; messages $(\pos_j, \vect{o}_j)$ from every neighbour $j \in N_i$, i.e.\ every drone within $R_{\mathrm{comm}}$; gains $k_c$, $k_l$; control period $\dt$; speed limit $v_{\max}$; the reference velocity $\dot{\vect{r}}$, forwarded from the leader through the graph to every drone (use $\vect{0}$ if unknown); if drone $i$ is pinned ($b_i = 1$), also the reference $\vect{r}$; optional extra velocity $\vel_i^{\mathrm{extra}}$ (avoidance, range keeping)}
\KwOut{the velocity command $\vel_i$ and the new position $\pos_i$}
\Fn{\ConsensusFormationStep{$i$}}{
  $\vel_i \gets \vect{0}$\;
  \ForEach{$j \in N_i$}{
    $\vect{d}_{ij} \gets \vect{o}_j - \vect{o}_i$\label{alg:ch23-formation:dij}\;
    $\vel_i \gets \vel_i + k_c\,(\pos_j - \pos_i - \vect{d}_{ij})$\label{alg:ch23-formation:sum}\tcp*{consensus on $\pos - \vect{o}$}
  }
  \If{$b_i = 1$}{
    $\vel_i \gets \vel_i + k_l\,(\vect{r} + \vect{o}_i - \pos_i)$\label{alg:ch23-formation:pin}\tcp*{pin the leader to the reference}
  }
  $\vel_i \gets \vel_i + \dot{\vect{r}} + \vel_i^{\mathrm{extra}}$\label{alg:ch23-formation:ff}\tcp*{feed-forward and extra terms}
  \If{$\norm{\vel_i} > v_{\max}$}{
    $\vel_i \gets v_{\max}\,\vel_i / \norm{\vel_i}$\label{alg:ch23-formation:sat}\;
  }
  $\pos_i \gets \pos_i + \dt\,\vel_i$\label{alg:ch23-formation:step}\tcp*{or hand $\vel_i$ to the autopilot}
  broadcast $(\pos_i, \vect{o}_i)$\label{alg:ch23-formation:send}\;
  \KwRet{$\vel_i$, $\pos_i$}\;
}
\end{algorithm}

Line~\ref{alg:ch23-formation:dij} derives the displacement from the offsets,
which is why the message carries $\vect{o}_j$ (or a slot number from which
$\vect{o}_j$ is looked up). Line~\ref{alg:ch23-formation:sum} is
\cref{eq:ch23-formation} with the gain $k_c$; with all offsets zero it is the
consensus protocol \cref{eq:ch23-protocol}, and over the whole swarm
lines~\ref{alg:ch23-formation:sum} and~\ref{alg:ch23-formation:step} are the
discrete-time protocol \cref{eq:ch23-discrete} with $\eps = k_c\dt$, so the
gain and the period must satisfy $k_c\dt < 1/d_{\max}$.
Line~\ref{alg:ch23-formation:pin} is the pinning term of
\cref{eq:ch23-pinned} applied to the template origin: the leader is pulled to
$\vect{r} + \vect{o}_i$, its own slot of a template whose origin is on the
reference. Line~\ref{alg:ch23-formation:ff} adds the reference velocity, which
is known to the leader from its path and is forwarded to the followers, and the
extra velocities of \cref{sec:ch23-range}: the safety term of the reactive
layer and the range-keeping term. Line~\ref{alg:ch23-formation:sat} enforces
the speed limit of \cref{ch:ch02} without changing the direction, and
line~\ref{alg:ch23-formation:step} is the Euler step, which on a real drone is
replaced by sending $\vel_i$ to the velocity controller of the autopilot.
Line~\ref{alg:ch23-formation:send} closes the loop: the position that drone
$i$ broadcasts now is what its neighbours use in their next period. The whole
period costs $\bigO{\abs{N_i}\,m}$ operations, one message in from each
neighbour and one message out.

\subsection{Second-order consensus for double-integrator drones}
\label{sec:ch23-second}

A quadrotor is not a first-order system that can set its velocity at will: to
a good approximation it is the double integrator $\ddot{\pos}_i = \ctrl_i$ of
\cref{ch:ch02}, whose input is an acceleration. Feeding the velocity of
\cref{eq:ch23-formation} through an inner velocity loop works when that loop
is fast; when the formation loop and the dynamics interact, one uses the
\textbf{second-order consensus}\index{consensus!second-order}\index{double integrator!consensus}
protocol \cite{ren2008distributed,olfatisaber2007consensus}
\begin{equation}
  \ctrl_i = \sum_{j \in N_i} \bigl[ k_p (\pos_j - \pos_i - \vect{d}_{ij})
  + k_v (\vel_j - \vel_i) \bigr] - k_d (\vel_i - \dot{\vect{r}})
  + b_i \bigl[ k_p (\vect{r} + \vect{o}_i - \pos_i) + k_v (\dot{\vect{r}} - \vel_i) \bigr].
  \label{eq:ch23-second}
\end{equation}
The first sum is the displacement term of \cref{eq:ch23-formation} used as a
spring. The second sum is consensus on the velocities: it damps the relative
motion of neighbours. The term $-k_d(\vel_i - \dot{\vect{r}})$ is absolute
damping towards the reference velocity (with $\dot{\vect{r}} = \vect{0}$ it
brings the swarm to rest); like the feed-forward of \cref{eq:ch23-pinned} it is
local only once every follower has been told $\dot{\vect{r}}$, forwarded from
the leader through the graph. The last bracket is the pinning of the leader.
Without any velocity term ($k_v = k_d = 0$) the drones are masses connected by
springs and oscillate for ever; with $k_v > 0$ or $k_d > 0$ every mode of the
formation decays (\cref{prop:ch23-second}). In the shifted variables, with
$\dot{\vect{r}}$ constant, the swarm obeys $\ddot{\vect{y}} =
-k_p\mat{L}\vect{y} - (k_v\mat{L} + k_d\mat{I})\dot{\vect{y}}
+ k_d\,\vect{1}\otimes\dot{\vect{r}}$ plus the pinning terms, which is again
the Laplacian acting on every coordinate. The forcing term
$k_d\,\vect{1}\otimes\dot{\vect{r}}$ is what makes the swarm travel with the
reference instead of coming to rest; it is the same on every drone, so it
moves the centroid and leaves the shape alone. \Cref{alg:ch23-second} is the corresponding
control period.

\begin{algorithm}[htb]
\caption{One control period of the second-order formation controller on drone $i$}
\label{alg:ch23-second}
\KwIn{own $\pos_i$, $\vel_i$, $\vect{o}_i$; messages $(\pos_j, \vel_j, \vect{o}_j)$ from $j \in N_i$; gains $k_p$, $k_v$, $k_d$; period $\dt$; the reference velocity $\dot{\vect{r}}$ on every drone (use $\vect{0}$ if unknown), the reference $\vect{r}$ if pinned}
\KwOut{acceleration command $\ctrl_i$; updated $\vel_i$, $\pos_i$}
\Fn{\ConsensusSecondOrderStep{$i$}}{
  $\ctrl_i \gets -k_d (\vel_i - \dot{\vect{r}})$\;
  \ForEach{$j \in N_i$}{
    $\ctrl_i \gets \ctrl_i + k_p (\pos_j - \pos_i - (\vect{o}_j - \vect{o}_i)) + k_v (\vel_j - \vel_i)$\label{alg:ch23-second:sum}\;
  }
  \If{$b_i = 1$}{$\ctrl_i \gets \ctrl_i + k_p (\vect{r} + \vect{o}_i - \pos_i) + k_v (\dot{\vect{r}} - \vel_i)$\;}
  $\vel_i \gets \vel_i + \dt\,\ctrl_i$; \ $\pos_i \gets \pos_i + \dt\,\vel_i$\label{alg:ch23-second:step}\tcp*{semi-implicit Euler}
  broadcast $(\pos_i, \vel_i, \vect{o}_i)$\;
  \KwRet{$\ctrl_i$, $\vel_i$, $\pos_i$}\;
}
\end{algorithm}

Line~\ref{alg:ch23-second:sum} needs the neighbours' velocities as well as
their positions, so the messages double in size, and
line~\ref{alg:ch23-second:step} updates the velocity before the position
(semi-implicit Euler), which keeps the oscillatory modes of the spring system
stable for the step sizes used in practice.

% ---------------------------------------------------------------------
\section{A worked example}
\label{sec:ch23-example}

\begin{example}[Four drones, one Laplacian, three controllers]
\label{ex:ch23-square}
Four drones hover at the planar positions
$\pos_1 = (0, 0)$, $\pos_2 = (2, 0)$, $\pos_3 = (1, 1.5)$ and
$\pos_4 = (3.5, 3)$ (metres) with a radio range $R_{\mathrm{comm}} = 3$~m.
The pairwise distances are $\ell_{12} = 2$, $\ell_{13} = \ell_{23} = 1.80$,
$\ell_{34} = 2.92$, $\ell_{24} = 3.35$ and $\ell_{14} = 4.61$, so the radius
graph has the four edges $\set{1,2}, \set{1,3}, \set{2,3}, \set{3,4}$
(\cref{fig:ch23-square}a, grey), the degrees $(2, 2, 3, 1)$ and
\[
  \mat{A} = \begin{pmatrix} 0&1&1&0\\ 1&0&1&0\\ 1&1&0&1\\ 0&0&1&0 \end{pmatrix}, \qquad
  \mat{L} = \mat{D} - \mat{A} = \begin{pmatrix} 2&-1&-1&0\\ -1&2&-1&0\\ -1&-1&3&-1\\ 0&0&-1&1 \end{pmatrix}.
\]
Every row sums to zero. The eigenvalues are $\lambda_1 = 0$, $\lambda_2 = 1$,
$\lambda_3 = 3$ and $\lambda_4 = 4$ (their sum, $8$, is the trace of
$\mat{L}$), with the eigenvectors $\vect{1}$, $(1, 1, 0, -2)\T$,
$(1, -1, 0, 0)\T$ and $(1, 1, -3, 1)\T$; multiply $\mat{L}$ by each to check.
The graph is connected ($\lambda_2 > 0$), its \textbf{Fiedler vector}\index{Fiedler vector}
$(1, 1, 0, -2)\T$ separates the leaf drone~4 from drones~1 and~2 across the
bottleneck edge $\set{3,4}$, and the \textbf{time constant}\index{consensus!time constant $1/\lambda_2$}
of consensus is $1/\lambda_2 = 1$~s. The bounds on the discrete step size are
$1/d_{\max} = 1/3$ and $2/\lambda_4 = 1/2$.

\emph{Consensus on the altitudes.} Let the drones agree on a common altitude
starting from $\vect{z}(0) = (10, 4, 7, 1)\T$~m, whose average is $5.5$~m.
\Cref{tab:ch23-consensus} lists the exact solution of $\dot{\vect{z}} =
-\mat{L}\vect{z}$ and the disagreement norm $\norm{\vect{\delta}(t)}$ next to
the bound $e^{-\lambda_2 t}\norm{\vect{\delta}(0)}$ of
\cref{cor:ch23-rate}. Drone~4, which hears only drone~3, is the slow one; after
$3$~s the bound leaves $5\%$ of the initial disagreement and after $5$~s
$0.7\%$, and the actual disagreement is smaller still because its components
along the faster modes $\lambda_3$ and $\lambda_4$ have long vanished.

\emph{Discrete time.} \Cref{tab:ch23-discrete} iterates
\cref{eq:ch23-discrete} from the same altitudes with $\eps = 0.25 < 1/3$ and
with $\eps = 0.6 > 1/2$. The first run is a slowed-down copy of the continuous
one (after the first step drone~3 lands exactly on the average of its three
neighbours, $(7.75 + 6.25 + 2.5)/3 = 5.5$, so its update term is zero and it
never moves again, while the others keep converging towards it); the second run
overshoots at every step, the mode $(1, 1, -3, 1)\T$ of $\lambda_4$ is
multiplied by $1 - 0.6 \cdot 4 = -1.4$ each time, and the altitudes fly apart.

\emph{Formation.} Now let the drones form the square with the offsets
$\vect{o}_1 = (0, 0)$, $\vect{o}_2 = (2, 0)$, $\vect{o}_3 = (2, 2)$,
$\vect{o}_4 = (0, 2)$ on the same fixed graph, which also serves as the
formation graph. The shifted variables are $\vect{y} = \pos - \vect{o} =
\bigl((0,0), (0,0), (-1, -0.5), (3.5, 1)\bigr)$; drones~1 and~2 already agree,
drone~4 is far off. Their average, the centroid shift, is
$\bar{\vect{y}} = (0.625, 0.125)$, so \cref{thm:ch23-formation} predicts the
final positions $\vect{o}_i + \bar{\vect{y}}$: $(0.625, 0.125)$,
$(2.625, 0.125)$, $(2.625, 2.125)$ and $(0.625, 2.125)$, the square translated by
$\bar{\vect{y}}$ (\cref{fig:ch23-square}a, circles). The initial formation error
over the four edges is $e_F(0) = 2.5$~m (the edge $\set{3,4}$ alone contributes
$\norm{(4.5, 1.5)}^2 = 22.5$~m$^2$ to the mean of squares) and the centred error
is $e_c(0) = 1.794$~m. \Cref{tab:ch23-formation} follows both errors under
\cref{eq:ch23-formation}; they decay with the time constant $1$~s of
$\lambda_2$, and $e_F$ drops below $1\%$ of its initial value after $4.06$~s.
All numbers in this example are produced by \code{worked\_example()} in
\code{code/ch23\_consensus.py}.
\end{example}

\begin{table}[tb]
  \centering\small
  \caption{Consensus on the altitudes of \cref{ex:ch23-square}: the exact
  solution of $\dot{\vect{z}} = -\mat{L}\vect{z}$, the disagreement norm and
  the bound $e^{-\lambda_2 t}\norm{\vect{\delta}(0)}$ with $\lambda_2 = 1$. The
  average $5.5$~m is conserved in every row.}
  \label{tab:ch23-consensus}
  \begin{tabular}{@{}rccccrr@{}}
  \toprule
  $t$ [s] & $z_1$ & $z_2$ & $z_3$ & $z_4$ & $\norm{\vect{\delta}(t)}$ & bound\\
  \midrule
  0   & 10.000 & 4.000 & 7.000 & 1.000 & 6.708 & 6.708\\
  0.5 & 7.315 & 5.976 & 5.703 & 3.006 & 3.127 & 4.069\\
  1   & 6.376 & 6.077 & 5.527 & 4.019 & 1.815 & 2.468\\
  2   & 5.778 & 5.763 & 5.501 & 4.958 & 0.663 & 0.908\\
  3   & 5.600 & 5.599 & 5.500 & 5.301 & 0.244 & 0.334\\
  5   & 5.513 & 5.513 & 5.500 & 5.473 & 0.033 & 0.045\\
  \bottomrule
  \end{tabular}
\end{table}

\begin{table}[tb]
  \centering\small
  \caption{Discrete-time consensus $\vect{z}_{k+1} = (\mat{I} - \eps\mat{L})\vect{z}_k$
  on \cref{ex:ch23-square}. Left: $\eps = 0.25$ satisfies
  $\eps < 1/d_{\max} = 1/3$ and converges. Right: $\eps = 0.6$ exceeds the
  exact bound $2/\lambda_4 = 0.5$ and diverges; the sign of the pattern
  $(1, 1, -3, 1)$ flips at every step.}
  \label{tab:ch23-discrete}
  \begin{tabular}{@{}r cccc c cccc@{}}
  \toprule
  & \multicolumn{4}{c}{$\eps = 0.25$} & & \multicolumn{4}{c}{$\eps = 0.6$}\\
  \cmidrule{2-5}\cmidrule{7-10}
  $k$ & $z_1$ & $z_2$ & $z_3$ & $z_4$ & & $z_1$ & $z_2$ & $z_3$ & $z_4$\\
  \midrule
  0 & 10.000 & 4.000 & 7.000 & 1.000 & & 10.000 & 4.000 & 7.000 & 1.000\\
  1 & 7.750 & 6.250 & 5.500 & 2.500 & & 4.600 & 9.400 & 3.400 & 4.600\\
  2 & 6.812 & 6.438 & 5.500 & 3.250 & & 6.760 & 2.920 & 8.440 & 3.880\\
  3 & 6.391 & 6.297 & 5.500 & 3.812 & & 5.464 & 8.536 & 1.384 & 6.616\\
  4 & 6.145 & 6.121 & 5.500 & 4.234 & & 4.859 & 2.402 & 11.262 & 3.477\\
  \bottomrule
  \end{tabular}
\end{table}

\begin{table}[tb]
  \centering\small
  \caption{Formation error of \cref{ex:ch23-square} under the
  displacement-based controller \cref{eq:ch23-formation} with $k_c = 1$: the
  edge RMS error $e_F$, the centred error $e_c$ and the bound
  $e_c(0)\,e^{-\lambda_2 t}$ of \cref{thm:ch23-formation}.}
  \label{tab:ch23-formation}
  \begin{tabular}{@{}rrrr@{}}
  \toprule
  $t$ [s] & $e_F(t)$ [m] & $e_c(t)$ [m] & $e_c(0)\,e^{-\lambda_2 t}$ [m]\\
  \midrule
  0 & 2.500 & 1.794 & 1.794\\
  1 & 0.542 & 0.541 & 0.660\\
  2 & 0.197 & 0.197 & 0.243\\
  3 & 0.072 & 0.072 & 0.089\\
  5 & 0.010 & 0.010 & 0.012\\
  \bottomrule
  \end{tabular}
\end{table}

\begin{figure}[tb]
  \centering
  \inputfigure{ch23/square}
  \caption{The four drones of \cref{ex:ch23-square} converging into a square
  under \cref{eq:ch23-formation}. (a) Trajectories from the initial positions
  (dots) along the fixed communication graph (grey) to the final square
  (circles), which is the template translated by the centroid shift
  $\bar{\vect{y}} = (0.625, 0.125)$; drones~1 and~2 move although they already
  agree with each other, because the template origin they end up on is the
  average of all four. (b) The formation error decays with the time constant
  $1/\lambda_2 = 1$~s; the dotted line is the bound of \cref{thm:ch23-formation}.}
  \label{fig:ch23-square}
\end{figure}

% ---------------------------------------------------------------------
\section{Properties: convergence, rate and step size}
\label{sec:ch23-properties}

\subsection{Average consensus on a connected graph}

\begin{theorem}[Average consensus, \textcite{olfatisaber2004consensus}]
\label{thm:ch23-consensus}
Let the communication graph be undirected and connected. Then for every
initial state $\vect{x}(0) \in \R^n$ the solution of
$\dot{\vect{x}} = -\mat{L}\vect{x}$ converges to the average of the initial
states, $\vect{x}(t) \to \bar{x}\vect{1}$ with $\bar{x} =
\frac{1}{n}\vect{1}\T\vect{x}(0)$, and the disagreement vector satisfies
\begin{equation}
  \norm{\vect{\delta}(t)} \le e^{-\lambda_2 t}\,\norm{\vect{\delta}(0)}
  \qquad\text{for all } t \ge 0.
  \label{eq:ch23-rate}
\end{equation}
\end{theorem}
\begin{proof}
By \cref{prop:ch23-laplacian}, $\mat{L}$ is symmetric, so it has an
orthonormal basis of eigenvectors $\vect{v}_1, \dots, \vect{v}_n$ with
eigenvalues $0 = \lambda_1 < \lambda_2 \le \dots \le \lambda_n$, where
$\vect{v}_1 = \vect{1}/\sqrt{n}$ and $\lambda_2 > 0$ because the graph is
connected. Expand the state in this basis, $\vect{x}(t) = \sum_k q_k(t)\,
\vect{v}_k$ with $q_k(t) = \vect{v}_k\T\vect{x}(t)$. Each coefficient obeys
$\dot{q}_k = \vect{v}_k\T\dot{\vect{x}} = -\vect{v}_k\T\mat{L}\vect{x} =
-\lambda_k \vect{v}_k\T\vect{x} = -\lambda_k q_k$, using
$\vect{v}_k\T\mat{L} = \lambda_k\vect{v}_k\T$ (symmetry). Hence
$q_k(t) = e^{-\lambda_k t} q_k(0)$ and
\begin{equation}
  \vect{x}(t) = q_1(0)\,\vect{v}_1 + \sum_{k=2}^{n} e^{-\lambda_k t}\, q_k(0)\,\vect{v}_k .
  \label{eq:ch23-modes}
\end{equation}
The first term is $\frac{\vect{1}\T\vect{x}(0)}{\sqrt{n}}\cdot\frac{\vect{1}}{\sqrt{n}}
= \bar{x}\vect{1}$, and it does not move: the average is conserved. The sum is
orthogonal to $\vect{1}$, so it is exactly the disagreement vector
$\vect{\delta}(t)$, and every one of its terms decays because $\lambda_k \ge
\lambda_2 > 0$. By orthonormality,
$\norm{\vect{\delta}(t)}^2 = \sum_{k \ge 2} e^{-2\lambda_k t} q_k(0)^2 \le
e^{-2\lambda_2 t} \sum_{k \ge 2} q_k(0)^2 = e^{-2\lambda_2 t}\norm{\vect{\delta}(0)}^2$.
\end{proof}

\begin{corollary}[Convergence rate]
\label{cor:ch23-rate}
Under the consensus protocol the disagreement decays at least as fast as
$e^{-\lambda_2 t}$, with the \textbf{time constant} $1/\lambda_2$, and no
faster bound holds for all initial states: if $\vect{\delta}(0)$ is a Fiedler
vector (an eigenvector of $\lambda_2$), \cref{eq:ch23-rate} is an equality.
With a gain $k_c$ in \cref{eq:ch23-protocol} the rate is $k_c\lambda_2$, and for
vector states \cref{eq:ch23-rate} holds coordinate by coordinate.
\end{corollary}

The corollary turns $\lambda_2$ into a design quantity: it says how long the
swarm needs to agree and how strongly it resists being pulled apart.
\Cref{fig:ch23-consensus} shows the effect for six drones in a line whose
initial states $(12, 3, 9, 1, 6, 11)$ have the average $7$. With a short radio
range each drone hears only its two neighbours, the graph is a path,
$\lambda_2 = 0.27$ and the time constant is $3.7$~s; with a longer range the
graph has nine edges, $\lambda_2 = 1.19$, and the time constant is $0.84$~s.
The disagreement falls below $1\%$ of its initial value after $6.3$~s on the
path and after $1.9$~s on the denser graph, and the slopes of the two decays in
the logarithmic plot are the two values of $\lambda_2$.

\begin{figure}[tb]
  \centering
  \inputfigure{ch23/consensus}
  \caption{Consensus of six drones on two radius graphs with different
  algebraic connectivity. (a) With $R_{\mathrm{comm}} = 2.1$~m the graph is a
  path ($\lambda_2 = 0.27$): the end drones are slow and the states creep
  towards the average $7$ (dotted). (b) With $R_{\mathrm{comm}} = 3.8$~m the
  graph has nine edges ($\lambda_2 = 1.19$) and agreement is reached four times
  faster. (c) The disagreement norms on a logarithmic scale with the bounds
  $e^{-\lambda_2 t}\norm{\vect{\delta}(0)}$ (dashed): the bound is tight in
  slope, and after the faster modes have died out the curves run parallel to
  it.}
  \label{fig:ch23-consensus}
\end{figure}

\begin{pitfall}[A disconnected graph agrees in clusters, not as a swarm]
Nothing in \cref{eq:ch23-protocol} tells a drone that the graph is
disconnected; each component runs its own consensus and settles on its own
average. Two pairs of drones $10$~m apart with $R_{\mathrm{comm}} = 2$~m and
the initial states $(1, 3, 10, 14)$ converge to $(2, 2, 12, 12)$, the two
cluster averages, and the disagreement norm settles at $10$ instead of decaying
to zero: every
drone sees its neighbours agreeing with it and is satisfied. The only local
symptom is that $\lambda_2 = 0$, which is why \cref{sec:ch23-range} monitors
it, and why the formation controller of a swarm must keep the graph connected
rather than hope that it is.
\end{pitfall}

\subsection{Discrete time and the step size}

\begin{theorem}[Discrete-time consensus]
\label{thm:ch23-discrete}
Let the graph be undirected and connected and let $0 < \eps < 2/\lambda_n$.
Then the iteration $\vect{x}_{m+1} = (\mat{I} - \eps\mat{L})\vect{x}_m$ of
\cref{eq:ch23-discrete}, whose time step is written $m$ here because $k$
indexes the modes in the proof, conserves the average and converges to it geometrically,
\begin{equation}
  \norm{\vect{\delta}_m} \le \rho^m \norm{\vect{\delta}_0}, \qquad
  \rho = \max\bigl(\abs{1 - \eps\lambda_2}, \abs{1 - \eps\lambda_n}\bigr) < 1 .
  \label{eq:ch23-discreterate}
\end{equation}
In particular the local condition $\eps < 1/d_{\max}$ of
\cref{eq:ch23-stepsize} suffices, because $\lambda_n \le 2 d_{\max}$.
\end{theorem}
\begin{proof}[Proof sketch]
$\mat{P} = \mat{I} - \eps\mat{L}$ has the eigenvectors of $\mat{L}$ with the
eigenvalues $\mu_k = 1 - \eps\lambda_k$; $\mu_1 = 1$ keeps the average, and
expanding the initial disagreement in the modes, $\vect{\delta}_0 =
\sum_{k \ge 2} c_k \vect{v}_k$, gives $\vect{\delta}_m = \mat{P}^m
\vect{\delta}_0 = \sum_{k \ge 2} c_k \mu_k^m \vect{v}_k$ after $m$ steps: mode
$k$ is damped by the factor $\mu_k^m$. All the $\mu_k$ with $k \ge 2$ lie
strictly inside $(-1, 1)$ exactly when
$0 < \eps\lambda_k < 2$ for $k \ge 2$, that is, when $\eps < 2/\lambda_n$
(the lower end is $\lambda_2 > 0$). The details, including the Gershgorin
step that gives $\lambda_n \le 2d_{\max}$, are \cref{exr:ch23-proof}.
\end{proof}

\begin{remark}[The best step size]
\label{rem:ch23-beststep}
$\rho$ is smallest when $1 - \eps\lambda_2 = -(1 - \eps\lambda_n)$, that is,
for $\eps^* = 2/(\lambda_2 + \lambda_n)$ \cite{xiao2004fast}. For the graph of
\cref{ex:ch23-square} this is $\eps^* = 0.4$ with $\rho = 0.6$ per step, while
the local bound $\eps = 1/3$ gives $\rho = 2/3$. The optimum needs the spectrum
of the whole graph, so in a swarm one settles for a safe margin below
$1/d_{\max}$, or for the weighted averages of \textcite{xiao2004fast}.
\end{remark}

\subsection{Directed graphs, switching topologies and delays}
\label{sec:ch23-directed}

On a directed graph, drone $i$ still runs \cref{eq:ch23-protocol} over the
drones it hears, and $\mat{L} = \mat{D} - \mat{A}$ with $d_i = \sum_j a_{ij}$
still satisfies $\mat{L}\vect{1} = \vect{0}$, but $\mat{L}$ is no longer
symmetric and the sum of the states is no longer conserved.
\textcite{ren2005consensus} proved the decisive result: the protocol reaches
consensus for every initial state if and only if the digraph contains a
\textbf{directed spanning tree}\index{spanning tree, directed}, a root drone
from which every other drone can be reached by following directed edges. Only
the root's information can propagate everywhere; two drones that hear nobody
can never agree. The agreed value is $\vect{w}\T\vect{x}(0)$, where
$\vect{w} \ge \vect{0}$ is the left eigenvector of $\mat{L}$ for the eigenvalue
$0$ scaled to $\vect{w}\T\vect{1} = 1$: a weighted average that gives weight
only to the drones that can reach all the others. If the digraph is
\textbf{balanced}, with in-degree equal to out-degree at every drone, then
$\vect{w} = \vect{1}/n$ and the average is recovered
\cite{olfatisaber2004consensus}. For \textbf{switching topologies}, when the
graph changes as the drones move or as packets are lost, consensus is still
reached if the union of the graphs over each of a sequence of bounded,
consecutive time intervals contains a spanning tree
\cite{ren2005consensus}; for undirected graphs this is the condition that the
union over bounded intervals is connected, first proved by
\textcite{jadbabaie2003coordination} for the alignment rule of
\cref{sec:ch23-flocking}. Finally, if every message arrives with the same
delay $\tau$, \textcite{olfatisaber2004consensus} showed that the protocol
still converges to the average if and only if $\tau < \pi/(2\lambda_n)$: a
dense graph (large $\lambda_n$) tolerates less delay, because it reacts faster
to information that is already stale.

\subsection{Leader--follower tracking and formation convergence}

\begin{proposition}[Leader--follower consensus]
\label{prop:ch23-leader}
Let the graph be undirected and connected, let at least one drone be pinned
and $k_l > 0$. Then $\mat{M} = \mat{L} + k_l\mat{B}$ is positive definite, and
under \cref{eq:ch23-pinned} every state converges to the reference: the
tracking error obeys $\norm{\vect{e}(t)} \le e^{-\mu_1 t}\norm{\vect{e}(0)}$,
where $\mu_1 > 0$ is the smallest eigenvalue of $\mat{M}$. Without the
feed-forward term the followers converge to a lagged reference,
$\vect{e}(t) \to -(\mat{M}\inv\vect{1})\otimes\dot{\vect{r}}$ for a constant
$\dot{\vect{r}}$.
\end{proposition}
\begin{proof}
$\mat{M}$ is symmetric and $\vect{x}\T\mat{M}\vect{x} = \sum_{\set{i,j}\in E}
(x_i - x_j)^2 + k_l\sum_{i:\,b_i = 1} x_i^2 \ge 0$. If the form is zero, the
first sum forces $\vect{x}$ to be constant on the connected graph and the
second forces that constant to be zero at a pinned drone, so $\vect{x} =
\vect{0}$: $\mat{M}$ is positive definite. The error
dynamics \cref{eq:ch23-pinnedmatrix} then give $\frac{d}{dt}\norm{\vect{e}}^2 =
-2\vect{e}\T\mat{M}\vect{e} \le -2\mu_1\norm{\vect{e}}^2$, and Gr\"onwall's
inequality yields the bound. Without feed-forward, $\dot{\vect{e}} =
-\mat{M}\vect{e} - \vect{1}\otimes\dot{\vect{r}}$ has the stable equilibrium
$-(\mat{M}\inv\vect{1})\otimes\dot{\vect{r}}$, since $\mat{M}$ is invertible.
\end{proof}

For the graph of \cref{ex:ch23-square} with drone~1 pinned and $k_l = 1$,
$\mu_1 = 0.18$, so a pinned formation tracks a step in the reference with the
time constant $5.6$~s, slower than the $1$~s of consensus among the drones:
the reference has to make its way through the graph. Without feed-forward at
$1$~m/s the lags are $4$~m for the leader and up to $6.67$~m for the tail
drone, and a formation error of $1.19$~m remains for ever
(\cref{exr:ch23-secondorder}). \Cref{fig:ch23-leader} shows a formation
forming behind a leader that follows a waypoint path with a switching radius
graph, with and without feed-forward.

\begin{figure}[tb]
  \centering
  \inputfigure{ch23/leader}
  \caption{Four drones forming a square of side $2$~m behind drone~1, which is
  pinned to a reference point (dashed) that moves along three waypoints at
  $1$~m/s; $R_{\mathrm{comm}} = 3.5$~m, $k_c = 1$, $k_l = 1.5$, $v_{\max} =
  3$~m/s. The polygons are snapshots of the formation. The formation error
  falls from $1.72$~m to $0.12$~m after $5$~s and to $0.0025$~m after $20$~s,
  and the switching radius graph stays connected throughout
  ($\lambda_2 \ge 1$). Without the feed-forward velocity $\dot{\vect{r}}$ the
  same run settles at a formation error of $1.00$~m: the followers trail the
  leader, and the leader trails the reference.}
  \label{fig:ch23-leader}
\end{figure}

\begin{theorem}[Formation convergence up to translation]
\label{thm:ch23-formation}
Let the communication graph be undirected and connected and let the
displacements come from offsets, $\vect{d}_{ij} = \vect{o}_j - \vect{o}_i$.
Under \cref{eq:ch23-formation} the positions converge to the formation
translated by the initial centroid shift,
\begin{equation}
  \pos_i(t) \to \vect{o}_i + \bar{\vect{y}}, \qquad
  \bar{\vect{y}} = \frac{1}{n}\sum_{j=1}^{n}\bigl(\pos_j(0) - \vect{o}_j\bigr),
  \label{eq:ch23-formlimit}
\end{equation}
the centred formation error satisfies $e_c(t) \le e^{-\lambda_2 t} e_c(0)$,
and when the formation graph is the communication graph the edge error is
squeezed between the two,
\begin{equation}
  \sqrt{\frac{n\lambda_2}{\abs{E}}}\; e_c(t) \;\le\; e_F(t) \;\le\; \sqrt{\frac{n\lambda_n}{\abs{E}}}\; e_c(t),
  \label{eq:ch23-sandwich}
\end{equation}
so $e_F(t) \to 0$ at the rate $\lambda_2$.
\end{theorem}
\begin{proof}
By \cref{eq:ch23-shift} the shifted variables $\vect{y}_i = \pos_i -
\vect{o}_i$ obey $\dot{\vect{y}} = -\mat{L}\vect{y}$ in every coordinate, and
\cref{thm:ch23-consensus} applied to each coordinate gives $\vect{y}_i(t) \to
\bar{\vect{y}}$, which is \cref{eq:ch23-formlimit}. The disagreement of
coordinate $c$ is $\vect{\delta}^{(c)} = \vect{y}^{(c)} - \bar{y}^{(c)}\vect{1}$
and $n\,e_c^2 = \sum_c \norm{\vect{\delta}^{(c)}}^2$, so \cref{eq:ch23-rate}
in every coordinate gives $e_c(t) \le e^{-\lambda_2 t}e_c(0)$. For the
sandwich, \cref{prop:ch23-laplacian}(ii) gives $\sum_{\set{i,j}\in E}
(y^{(c)}_i - y^{(c)}_j)^2 = \vect{\delta}^{(c)\mathsf{T}}\mat{L}\vect{\delta}^{(c)}$,
which lies between $\lambda_2\norm{\vect{\delta}^{(c)}}^2$ and
$\lambda_n\norm{\vect{\delta}^{(c)}}^2$ because $\vect{\delta}^{(c)} \perp
\vect{1}$; summing over the coordinates and dividing by $\abs{E}$ gives
$\lambda_2 n e_c^2 \le \abs{E}\, e_F^2 \le \lambda_n n e_c^2$.
\end{proof}

In \cref{ex:ch23-square}, $n = \abs{E} = 4$, so $e_c \le e_F \le 2e_c$; the
table shows $e_F \approx e_c$ from $t = 1$~s on, because by then only the
$\lambda_2$ mode is left and the lower bound is attained. The theorem also
answers the question where the formation ends up: exactly where the initial
positions put the template's centroid. A pinned leader overrides this, and
\cref{prop:ch23-leader} applied to $\vect{y}$ shows that the formation then
converges to $\vect{o}_i + \vect{r}$.

\begin{proposition}[Second-order formation control]
\label{prop:ch23-second}
Let the graph be undirected and connected and let the drones be double
integrators under \cref{eq:ch23-second} without pinning, with a constant
reference velocity $\dot{\vect{r}}$. In the eigenbasis of $\mat{L}$ the
shifted dynamics decouple into the disagreement modes
\begin{equation}
  \ddot{q}_k + (k_v\lambda_k + k_d)\,\dot{q}_k + k_p\lambda_k\, q_k = 0,
  \qquad k = 2, \dots, n .
  \label{eq:ch23-modes2}
\end{equation}
The centroid mode $k = 1$ ($\lambda_1 = 0$, $\vect{v}_1 = \vect{1}/\sqrt{n}$)
carries the whole forcing and obeys $\ddot{q}_1 + k_d\,\dot{q}_1 =
k_d\sqrt{n}\,\dot{r}$ instead, where $\dot{r}$ is the component of
$\dot{\vect{r}}$ in the coordinate considered, so the centroid velocity
converges to $\dot{r}$ when $k_d > 0$ and is conserved when $k_d = 0$.
Every mode with $\lambda_k > 0$ decays if and only if $k_p > 0$ and $k_v\lambda_k
+ k_d > 0$, so the formation is reached for any $k_p > 0$ with $k_v > 0$ or
$k_d > 0$; with $k_v = k_d = 0$ every mode with $\lambda_k > 0$ oscillates for
ever at the angular frequency $\sqrt{k_p\lambda_k}$, while the centroid keeps
whatever velocity it started with.
\end{proposition}
\begin{proof}[Proof sketch]
$\ddot{\vect{y}} = -k_p\mat{L}\vect{y} - (k_v\mat{L} + k_d\mat{I})\dot{\vect{y}}
+ k_d\,\vect{1}\otimes\dot{\vect{r}}$
is diagonalised by the eigenvectors of $\mat{L}$, which are shared by
$\mat{I}$; the forcing is a multiple of $\vect{1} = \sqrt{n}\,\vect{v}_1$, so
it appears in the mode $k = 1$ alone and every mode $k \ge 2$ is homogeneous.
Each such mode is a damped oscillator with the characteristic polynomial
$s^2 + (k_v\lambda_k + k_d)s + k_p\lambda_k$, which has both roots in the open
left half-plane exactly when both coefficients are positive
(\cref{exr:ch23-secondorder}).
\end{proof}

The mode analysis is also the tuning rule. The damping ratio of mode $k$ is
$\zeta_k = (k_v\lambda_k + k_d)/(2\sqrt{k_p\lambda_k})$; for the graph of
\cref{ex:ch23-square} with $k_p = 1$ the undamped frequencies are $1$,
$\sqrt3$ and $2$~rad/s, and $k_v = 1$, $k_d = 0.5$ give $\zeta = 0.75$,
$1.01$ and $1.13$, close to critical for all three. In the simulation of
\code{code/ch23\_consensus.py} these gains bring the formation error below
$10^{-3}$~m in $25$~s with the drones at rest, while $k_v = k_d = 0$ leaves an
error above $0.1$~m that never decays.

% ---------------------------------------------------------------------
\section{Variants and extensions}
\label{sec:ch23-variants}

\subsection{Distance-based formations and rigidity}
\label{sec:ch23-distance}

The controller of \cref{sec:ch23-formation} needs every drone to express
$\pos_j - \pos_i$ in a frame whose orientation all drones share: a compass, a
GPS heading or the map frame of the global planner. Where no such frame
exists, \textbf{distance-based formation control}\index{formation control!distance-based}
prescribes only the lengths $d_{ij} = \norm{\vect{d}_{ij}}$ and lets drone $i$
move along the line to each neighbour, $\dot{\pos}_i = \sum_j
(\norm{\pos_j - \pos_i}^2 - d_{ij}^2)(\pos_j - \pos_i)$, which every drone can
compute in its own body frame \cite{krick2009stabilisation,oh2015survey}. The
price is high. Distances define a shape only up to rotation and reflection,
and only if the formation graph is \textbf{rigid}\index{rigidity}: a square
with four edges can fold into a rhombus, so a diagonal must be added, and in
general the graph must be infinitesimally rigid, a condition on the rank of a
rigidity matrix \cite{anderson2008rigid}. Moreover the dynamics are nonlinear,
convergence is guaranteed only from nearby configurations, and mirrored
formations are equilibria too. \Cref{tab:ch23-families} places the three
families of \textcite{oh2015survey} side by side. Since every drone in this
book shares the frame of the global planner, displacement-based control is the
right choice: it is linear, it converges from everywhere on any connected
graph, and the \cbs reference enters it through one pinned drone.

\begin{table}[tb]
  \centering\small
  \caption{The three families of formation control \cite{oh2015survey}. The
  displacement-based family is the one used in this book.}
  \label{tab:ch23-families}
  \begin{tabular}{@{}lL{3.0cm}L{2.6cm}L{2.6cm}L{2.6cm}@{}}
  \toprule
  Family & Each drone measures & Frame needed & Shape fixed up to & Graph condition\\
  \midrule
  Position-based & its own absolute position & global frame and heading & nothing (unique) & none (interaction optional)\\
  Displacement-based & relative positions $\pos_j - \pos_i$ & common heading only & translation & connected\\
  Distance-based & distances $\norm{\pos_j - \pos_i}$ & none (body frames) & translation, rotation, reflection & rigid\\
  \bottomrule
  \end{tabular}
\end{table}

\subsection{Flocking}
\label{sec:ch23-flocking}

A formation fixes every relative position. A \textbf{flock}\index{flocking}
fixes none of them and still moves as a group: the drones keep a comfortable
distance from each other, move in the same direction and follow a common
goal. \textcite{reynolds1987flocks} produced this behaviour in computer
graphics with three rules for his \emph{boids}\index{boids}: separation
(steer away from crowded neighbours), alignment (steer towards the average
heading of the neighbours) and cohesion (steer towards their average
position). \textcite{olfatisaber2006flocking} turned the rules into a
controller with a proof. Each double-integrator drone applies the acceleration
\begin{equation}
  \ctrl_i = \underbrace{\sum_{j \in N_i} \phi\bigl(\norm{\pos_j - \pos_i}\bigr)\,
  \frac{\pos_j - \pos_i}{\norm{\pos_j - \pos_i}}}_{\text{gradient: separation and cohesion}}
  + \underbrace{\sum_{j \in N_i} a_{ij}(\pos)\,(\vel_j - \vel_i)}_{\text{velocity alignment}}
  + \underbrace{\bigl(-c_1(\pos_i - \pos_\gamma) - c_2(\vel_i - \vel_\gamma)\bigr)}_{\text{navigation feedback}},
  \label{eq:ch23-flock}
\end{equation}
with the three terms of \cref{fig:ch23-flocking}. The \emph{gradient term} is
minus the gradient of a smooth pairwise potential (\cref{fig:ch23-potentials}
is its relative): the action function $\phi(\ell)$ is negative below the
desired distance $d$ (repulsion), positive between $d$ and the interaction
range $r$ (attraction) and exactly zero beyond $r$, where it is faded out by a
bump function so that neighbours can enter and leave the range without a jump
in the force; Olfati-Saber evaluates the distances through a smoothed norm
$\norm{\vect{z}}_\sigma = (\sqrt{1 + \sigma\norm{\vect{z}}^2} - 1)/\sigma$ with a
small $\sigma > 0$, so that the potential is differentiable everywhere; he calls
this parameter $\epsilon$, which in this chapter is the step size of
\cref{eq:ch23-discrete}. The \emph{alignment term} is
consensus on the velocities with the same distance-dependent weights
$a_{ij}(\pos) \in [0, 1]$, so \cref{thm:ch23-consensus} explains why the
velocities become equal. The \emph{navigation feedback} is a PD term towards a
virtual leader (the $\gamma$-agent) at $\pos_\gamma$ with velocity
$\vel_\gamma$, for us the \cbs reference of the group. It is not an
afterthought: without it the flock fragments into small groups that drift
apart, which is one of the main results of \textcite{olfatisaber2006flocking}
and a warning against the pure boids rules for a swarm that must stay
together. Obstacles enter through virtual $\beta$-agents on their boundary,
exactly the mechanism of the repulsive potentials of \cref{ch:ch15}.

\begin{figure}[tb]
  \centering
  \inputfigure{ch23/flocking}
  \caption{Left: the three terms of \cref{eq:ch23-flock} acting on drone $i$;
  the neighbours within the range $r$ produce the gradient term (repel when
  too close, attract when too far) and the alignment term (match their mean
  velocity), the target produces the navigation term. Right: the action
  function of the gradient term, $\phi(\ell) < 0$ below the desired distance
  $d$, $\phi(\ell) > 0$ between $d$ and $r$, and zero beyond $r$, faded out
  smoothly by the bump function.}
  \label{fig:ch23-flocking}
\end{figure}

The connections to the rest of the book are direct. The gradient term is an
artificial potential field between agents in the sense of \cref{ch:ch15}
(Khatib's repulsion \cite{khatib1986realtime} with a cohesive part added), and
it inherits the weaknesses of potential fields: it is soft, it can be
overpowered, and it makes no promise about the minimum distance. In the hybrid
architecture (\cref{ch:ch24}) the separation is therefore delegated to the
reactive layer: \orca (\cref{ch:ch13}) takes the velocity computed by the
consensus or formation terms as its \emph{preferred velocity} and returns the
closest velocity that is collision-free for the \orca horizon $\ttc$ of
\cref{ch:ch13}, a hard
guarantee that a potential cannot give. Consensus decides where the swarm
wants to go, \orca decides what it may do this instant.

\begin{historynote}
Reynolds' boids (1987) came from computer animation and had no proof.
\textcite{vicsek1995novel} studied the alignment rule alone as a physics model
and observed a phase transition to ordered motion; \textcite{jadbabaie2003coordination}
gave the first convergence proof for it, using products of stochastic matrices
on switching graphs, which is the discrete-time consensus of
\cref{thm:ch23-discrete} in disguise. \textcite{olfatisaber2004consensus} set
up the Laplacian framework with $\lambda_2$ as the rate and treated delays and
switching, \textcite{ren2005consensus} settled the directed case with the
spanning-tree condition, \textcite{olfatisaber2006flocking} and
\textcite{tanner2007flocking} proved flocking, and the algebraic connectivity
itself goes back to \textcite{fiedler1973algebraic}.
\end{historynote}

\subsection{Communication-range constraints and connectivity maintenance}
\label{sec:ch23-range}

Every theorem above assumes a connected graph, and the radius graph of a
moving swarm is connected only as long as the drones keep it so.
\textbf{Connectivity maintenance}\index{connectivity maintenance} selects a
set of edges $E_{\mathrm{keep}}$ that spans the graph, typically the edges of
the formation graph, and keeps each of them shorter than $R_{\mathrm{comm}}$.
There are two ways to do this, and a quantity to watch.

\paragraph{As a hard constraint.}
The condition $\norm{\pos_i - \pos_j} \le R_{\mathrm{comm}}$ is convex in the
pair of positions (a disc of relative positions), so it can be handed to an
optimisation-based controller. In the MPC of \cref{ch:ch21} it becomes one
constraint per kept edge and per step of the horizon, either as a second-order
cone or, for a quadratic program, replaced by a polygon inscribed in the disc:
$\vect{n}_q\T(\pos_i - \pos_j) \le R_{\mathrm{comm}}\cos(\pi/Q)$ for $Q$ unit
directions $\vect{n}_q$, a slightly conservative inner approximation that
stays linear (\cref{exr:ch23-range}). The MILP formulations of \cref{ch:ch22}
use the same linear inequalities. The constraint is exact and can be combined
with the collision constraints, but it needs a planner that solves a program
at every period; for the distributed controller of \cref{alg:ch23-formation}
one uses the second way.

\paragraph{As a potential.}
Add to the displacement potential a \textbf{range-keeping barrier}\index{range-keeping barrier}
on each kept edge that is zero while the edge is comfortably short and grows
without bound as the length approaches $R_{\mathrm{comm}}$
\cite{ji2007distributed,zavlanos2011graph}:
\begin{equation}
  V_{\mathrm{range}}(\ell) = \begin{cases}
    0 & \ell \le \ell_{\mathrm{act}},\\[2pt]
    \dfrac{(\ell - \ell_{\mathrm{act}})^2}{R_{\mathrm{comm}} - \ell} & \ell_{\mathrm{act}} < \ell < R_{\mathrm{comm}},
  \end{cases}
  \label{eq:ch23-barrier}
\end{equation}
and let both drones of the edge descend its gradient (\cref{fig:ch23-potentials}).
Because the barrier is infinite at the range boundary, a finite disturbance
cannot stretch the edge beyond it in the continuous-time ideal; in discrete
time the barrier must be paired with the velocity limit so that one period
cannot jump across the last few centimetres, and the activation distance
$\ell_{\mathrm{act}}$ (say $0.9\,R_{\mathrm{comm}}$) leaves the formation
controller alone until an edge is actually at risk. The safety repulsion of
the reactive layer is the same construction at the other end of the scale,
infinite at the contact distance. The velocity of every term is added on
line~\ref{alg:ch23-formation:ff} of \cref{alg:ch23-formation}.

\begin{figure}[tb]
  \centering
  \inputfigure{ch23/potentials}
  \caption{The pairwise potentials acting on one edge as functions of its
  length. The displacement term (dashed) is a quadratic bowl around the desired
  length; the safety repulsion (green) is Khatib's potential, active below
  $\rho_0$ and infinite at contact; the range-keeping barrier (orange) is zero
  up to $\ell_{\mathrm{act}}$ and infinite at $R_{\mathrm{comm}}$. Formation
  control lives in the bowl, and the two barriers keep the edge out of the
  forbidden ends.}
  \label{fig:ch23-potentials}
\end{figure}

\paragraph{What to monitor.}
The barrier protects the edges of $E_{\mathrm{keep}}$; $\lambda_2$ of the
current radius graph tells whether the whole graph is still connected and
how robustly. It is the natural \textbf{connectivity measure}\index{algebraic connectivity $\lambda_2$!as a monitor}
of a swarm: zero means split, small means one edge away from splitting, and it
is also the convergence rate the formation controller currently enjoys. A
ground station with all positions computes it with an eigensolver in
$\bigO{n^3}$. On board, a drone can estimate it by power iteration on the
matrix $\mat{W} = s\mat{I} - \mat{L} - \frac{s}{n}\vect{1}\vect{1}\T$ with
$s = 2d_{\max}$ (a different matrix from the pinning matrix $\mat{B}$ of
\cref{eq:ch23-pinned}), whose largest eigenvalue is $s - \lambda_2$; each iteration
needs only products with $\mat{L}$ (one message exchange) and the average of
the iterate (a consensus), which is the basis of the distributed estimators of
\textcite{yang2010decentralized}. The function
\code{fiedler\_power\_iteration()} in \code{code/ch23\_consensus.py} is the
centralised version and agrees with NumPy's eigenvalues to $10^{-6}$ in the
self-test. Together with the lengths of the kept edges, $\lambda_2$ is the
connectivity constraint that \cref{ch:ch24} monitors and \cref{ch:ch25}
reports.

% ---------------------------------------------------------------------
\section{Implementation notes}
\label{sec:ch23-implementation}

\paragraph{Step size.}
The gain and the control period enter only as the product $\eps = k_c\dt$, and
\cref{thm:ch23-discrete} needs $\eps < 1/d_{\max}$ with margin. At $20$~Hz
($\dt = 0.05$~s) with $k_c = 1$~s$^{-1}$ and at most five neighbours,
$\eps = 0.05$ is far below $0.2$; the bound bites when the gain is raised for
a faster response, or when a slow link forces a long period. The range barrier
and the safety repulsion are stiffer than the consensus term near their
poles, so the velocity limit $v_{\max}$ on line~\ref{alg:ch23-formation:sat}
is what keeps the Euler step from jumping over a barrier.

\begin{pitfall}[An unstable step size looks like a bug in the messages]
With $\eps$ above $2/\lambda_n$ the swarm does not merely converge slowly: it
diverges, and it does so in the highest-frequency pattern of the graph, so
the symptom is neighbouring drones lurching in opposite directions at every
period (right half of \cref{tab:ch23-discrete}). Because $\lambda_n$ depends on
the graph, a gain that was fine in a sparse test can fail when the drones
cluster and the degrees grow. Also beware of the boundary: a $4$-cycle has
$\lambda_4 = 4 = 2d_{\max}$, so $\eps = 1/d_{\max} = 0.5$ exactly gives the
eigenvalue $-1$ and the pattern $(+,-,+,-)$ flips for ever without decaying.
Use a strict margin, and log $\max_i d_i$ together with the control period.
\end{pitfall}

\paragraph{Asynchronous updates and packet loss.}
Real drones do not update in lockstep. Using the latest position received from
each neighbour is a bounded delay and a lost packet is a switching topology, so
\cref{sec:ch23-directed} already covers both. What it does not cover is an
\emph{asymmetric} loss: if $i$ hears $j$ but $j$ misses $i$, the graph is
momentarily directed and unbalanced, and the agreed value drifts. A pinned
leader corrects that for a formation; a swarm that averages sensor readings
needs a protocol that retransmits or that exchanges mass explicitly.

\paragraph{Switching topologies and hysteresis.}
The radius graph switches whenever a pair crosses $R_{\mathrm{comm}}$, and a
pair hovering at exactly that distance would switch at every period. Use
hysteresis: add an edge when the pair comes within $0.9\,R_{\mathrm{comm}}$
and drop it only when it exceeds $R_{\mathrm{comm}}$. Keep the formation graph
$G_F$ separate from the communication graph: the controller uses whichever
edges exist, the formation error is always measured on $G_F$, and the range
barrier guards $E_{\mathrm{keep}} \subseteq E_F$.

\paragraph{The code.}
\code{code/ch23\_consensus.py} implements the chapter: the graph tools of
\cref{lst:ch23-graph}, the exact and RK4 continuous-time consensus, the
discrete protocol, the pinned matrix, the formation controllers of
\cref{lst:ch23-velocity} (first order) and \code{formation\_acceleration()}
(second order), the error metrics of \cref{lst:ch23-error}, the barrier and
repulsion velocities, a first-order simulator with a switching radius graph,
a second-order simulator, the flocking controller and the worked example.
Positions are the rows of an $(n, m)$ array, and the desired displacements are
never stored: the functions receive the offsets and compute
$\vect{d}_{ij} = \vect{o}_j - \vect{o}_i$ themselves. The self-test checks the
Laplacian identities and the power-iteration monitor against NumPy on random
radius graphs, the spectrum $\set{0, 1, 3, 4}$ and the numbers of
\cref{ex:ch23-square}, convergence to the average on a connected graph and to
cluster averages on a disconnected one, the step-size bounds on the $4$-cycle
and on random graphs, leader--follower convergence and its failure when the
pinned part is cut off, the formation-error identities and the convergence of
the first- and second-order controllers, the range barrier under a stretching
disturbance, and velocity alignment without collisions in the flock. It runs in
a few seconds.

\begin{lstlisting}[caption={The communication graph from positions and a radius, its Laplacian, its edge list and $\lambda_2$ (from \code{code/ch23\_consensus.py}).},label={lst:ch23-graph}]
def radius_graph(positions, r_comm):
    """Return the 0/1 adjacency matrix of the communication graph.

    Drones i != j are neighbours if and only if ||p_i - p_j|| <= r_comm.
    """
    P = np.asarray(positions, dtype=float)
    dist = np.linalg.norm(P[:, None, :] - P[None, :, :], axis=2)
    A = (dist <= r_comm).astype(float)
    np.fill_diagonal(A, 0.0)
    return A


def laplacian(A):
    """Graph Laplacian L = D - A; every row of L sums to zero."""
    A = np.asarray(A, dtype=float)
    return np.diag(A.sum(axis=1)) - A


def edges_of(A):
    """Undirected edges (i, j) with i < j and a_ij > 0."""
    A = np.asarray(A)
    n = A.shape[0]
    return [(i, j) for i in range(n) for j in range(i + 1, n) if A[i, j] > 0]


def algebraic_connectivity(L):
    """lambda_2(L): the second-smallest eigenvalue, zero iff disconnected."""
    return float(np.linalg.eigvalsh(np.asarray(L, dtype=float))[1])
\end{lstlisting}

\begin{lstlisting}[caption={The first-order displacement-based formation controller with leader pinning and feed-forward: \cref{alg:ch23-formation} for the whole swarm at once, using $\pos_j - \pos_i - \vect{d}_{ij} = \vect{y}_j - \vect{y}_i$ (from \code{code/ch23\_consensus.py}).},label={lst:ch23-velocity}]
def formation_velocity(P, offsets, A, gain=1.0, ref=None, ref_vel=None,
                       pinned=(), gain_ref=1.0):
    """First-order displacement-based formation controller (velocities).

    v_i = gain * sum_j a_ij (p_j - p_i - d_ij)
          + b_i gain_ref (r + o_i - p_i)     pinning of the leader(s)
          + rdot                              feed-forward if known
    Because p_j - p_i - d_ij = y_j - y_i with y = p - o, the sum is -L y.
    """
    P = np.asarray(P, dtype=float)
    O = np.asarray(offsets, dtype=float)
    V = -gain * (laplacian(A) @ (P - O))
    if ref is not None:
        r = np.asarray(ref, dtype=float)
        for i in pinned:
            V[i] += gain_ref * (r + O[i] - P[i])
    if ref_vel is not None:
        V += np.asarray(ref_vel, dtype=float)
    return V
\end{lstlisting}

\begin{lstlisting}[caption={The formation error of \cref{def:ch23-error} over the edges of the formation graph, and its centred variant (from \code{code/ch23\_consensus.py}).},label={lst:ch23-error}]
def formation_error(P, offsets, A_form):
    """RMS of ||(p_j - p_i) - d_ij|| over the edges of the formation graph."""
    Y = np.asarray(P, dtype=float) - np.asarray(offsets, dtype=float)
    edges = edges_of(A_form)
    if not edges:
        return float("nan")
    sq = [float(np.sum((Y[j] - Y[i]) ** 2)) for i, j in edges]
    return float(np.sqrt(np.mean(sq)))


def formation_error_centred(P, offsets):
    """RMS deviation from the template after removing the best translation."""
    Y = np.asarray(P, dtype=float) - np.asarray(offsets, dtype=float)
    Y = Y - Y.mean(axis=0)
    return float(np.sqrt(np.mean(np.sum(Y ** 2, axis=1))))
\end{lstlisting}

% ---------------------------------------------------------------------
\section{Where this fits in the drone system}
\label{sec:ch23-drone}

\begin{dronebox}
Consensus is the \textbf{low-level formation controller} of the hybrid
architecture of \cref{ch:ch24}, and the communication radius is the last of
its constraints. It sits between two layers and talks to both.

\textbf{From the global planner.} \cbs or \ecbs (\cref{ch:ch09,ch:ch10}) plan
one path for the formation as a whole, which MPC (\cref{ch:ch21}) may smooth;
sampled at the current time it is the reference $\vect{r}(t)$,
$\dot{\vect{r}}(t)$ of \cref{alg:ch23-formation}. One drone is pinned to
$\vect{r}$ and $\dot{\vect{r}}$ is forwarded through the graph to the others,
which need nothing else but their neighbours' positions and the offsets
$\vect{o}_i$ --- the shape of the mission (a line for a survey, a square for a
camera rig), with the graph connected and $\lambda_2$ as large as the radio
range allows.

\textbf{To and from the reactive layer.} The velocity of
\cref{alg:ch23-formation} is the \emph{preferred velocity} handed to \orca
(\cref{ch:ch13}) or DWA (\cref{ch:ch14}) for the unexpected obstacles that
the prediction layer (\cref{ch:ch18,ch:ch19,ch:ch20}) reports. The safety
layer may change it, and when it does the formation breaks temporarily.
That is intended: separation from an intruder has priority over the shape.
The consensus term keeps acting on the drones that were not pushed, so they
move to accommodate the one that was, and it pulls the formation back
together as soon as the intruder has passed. \Cref{fig:ch23-avoidance}
quantifies this on the square of \cref{ex:ch23-square}.

\textbf{Monitored constraints.} Two quantities are logged at every period and
reported by the evaluation of \cref{ch:ch25}: the formation error $e_F$ of
\cref{def:ch23-error} and the connectivity, as $\lambda_2$ of the current radius
graph together with the lengths of the kept edges and the count of
communication violations. The range barrier of \cref{sec:ch23-range} acts on
the same edges, and a drop of $\lambda_2$ towards zero, or a kept edge beyond
$\ell_{\mathrm{act}}$, is a trigger for the replanning layer (\dstarlite,
\astar; \cref{ch:ch05,ch:ch04}) to slow the reference down or to route the
formation around the obstacle as a group. This is the Week~11 material of the
study plan (\cref{ch:appA}); its coding exercise is \cref{exr:ch23-coding}.
\end{dronebox}

\begin{figure}[tb]
  \centering
  \inputfigure{ch23/avoidance}
  \caption{The square formation of \cref{ex:ch23-square} following a straight
  reference at $1$~m/s while an intruder crosses its path from above at
  $1$~m/s; $R_{\mathrm{comm}} = 3.2$~m, $\ell_{\mathrm{act}} = 2.95$~m, safety
  distance $d_{\mathrm{safe}} = 0.8$~m, repulsion within $2$~m. The shaded
  band marks the interval $[6.9, 11.7]$~s during which the repulsion is
  active. Top: the formation error rises to a peak of $0.36$~m at
  $t = 10.5$~s and is back below $0.05$~m at $t = 13.6$~s; the distance to the
  intruder is drawn divided by ten so that both curves fit one axis, and its
  smallest value is a closest approach of $0.91$~m. Bottom: the longest edge reaches
  $3.04$~m, held below $R_{\mathrm{comm}}$ by the range barrier, and the
  shortest edge stays at $1.56$~m; the graph never loses an edge
  ($\lambda_2 = 4$ throughout) and there are no communication violations.}
  \label{fig:ch23-avoidance}
\end{figure}

\begin{pitfall}[Avoidance breaks the formation, and must be allowed to]
A tempting fix for a high formation error during an avoidance manoeuvre is to
raise the formation gain until the shape holds. This turns the formation
controller into an adversary of the safety layer: the consensus term pulls the
avoiding drone back towards the intruder, the safety term pushes it out again,
and the drone oscillates at the edge of the safety distance. Treat the
formation error during avoidance as a cost to be reported (\cref{ch:ch25}),
not a constraint to be enforced, keep the gains such that the safety term
dominates near contact, and judge the controller by the recovery time
afterwards, $3$~s in \cref{fig:ch23-avoidance}. What must be enforced is the
other constraint: the range barrier, which keeps the stretched edges below
$R_{\mathrm{comm}}$ so that the swarm stays connected while it is out of
shape.
\end{pitfall}

% ---------------------------------------------------------------------
\section{Summary}
\begin{summary}
\begin{itemize}
  \item The communication graph of a swarm is the radius graph of the drone
    positions; its Laplacian $\mat{L} = \mat{D} - \mat{A}$ is symmetric,
    positive semidefinite, has $\mat{L}\vect{1} = \vect{0}$, and its second
    eigenvalue $\lambda_2$ is positive exactly when the graph is connected.
  \item The consensus protocol $\dot{\vect{x}} = -\mat{L}\vect{x}$ is local,
    leaderless and conserves the sum. On a connected undirected graph it
    converges to the average of the initial states, and the disagreement decays
    like $e^{-\lambda_2 t}$ (\cref{thm:ch23-consensus}).
  \item In discrete time, $\vect{x}_{k+1} = (\mat{I} - \eps\mat{L})\vect{x}_k$
    converges for $\eps < 2/\lambda_n$, and $\eps < 1/d_{\max}$ is the
    sufficient condition every drone can check (\cref{thm:ch23-discrete}).
  \item On a directed graph consensus needs a directed spanning tree and
    reaches a weighted average; on switching graphs it needs the union over
    bounded windows to be connected.
  \item Pinning one drone to a reference makes the whole graph follow it;
    the feed-forward $\dot{\vect{r}}$ removes the lag.
  \item Displacement-based formation control is consensus on the shifted
    positions $\pos_i - \vect{o}_i$; it converges to the formation up to a
    translation (\cref{thm:ch23-formation}). The formation error $e_F$ is the
    RMS violation of the desired displacements over the formation graph.
  \item Double-integrator drones need velocity consensus or absolute damping,
    otherwise the formation oscillates for ever.
  \item Flocking adds a gradient term for separation and cohesion, velocity
    alignment and navigation feedback; in the drone system the separation is
    delegated to \orca.
  \item Connectivity is kept by a hard constraint in MPC or by a barrier
    potential at $R_{\mathrm{comm}}$, and monitored through $\lambda_2$ and the
    edge lengths; avoidance breaks the formation temporarily, and the
    controller is judged by its recovery.
\end{itemize}
\end{summary}

\paragraph{Further reading.}
The history note above traces the original papers; this paragraph adds what it
does not name. The survey of \textcite{olfatisaber2007consensus} is the
standard entry point and covers everything in \cref{sec:ch23-properties} with
the Perron-matrix view of discrete time. The book of
\textcite{ren2008distributed} and the tutorial of \textcite{ren2007information}
treat leader--follower and second-order consensus and formation control for
vehicles; \textcite{mesbahi2010graph} and \textcite{bullo2009distributed} are
the graph-theoretic textbooks. \textcite{oh2015survey} organises formation
control into the three families of \cref{tab:ch23-families}, with
\textcite{anderson2008rigid} and \textcite{krick2009stabilisation} for the
rigidity theory of the distance-based family. Connectivity maintenance is
\textcite{ji2007distributed} and \textcite{zavlanos2011graph}, the distributed
estimation of $\lambda_2$ is \textcite{yang2010decentralized}, and the optimal
step size is \textcite{xiao2004fast}.

% ---------------------------------------------------------------------
\section{Exercises}
\label{sec:ch23-exercises}

\begin{exercise}[\difficulty{1}]
\label{exr:ch23-laplacian}
Four drones communicate over the edges $\set{1,2}$, $\set{2,3}$, $\set{3,4}$,
$\set{4,1}$ and $\set{1,3}$ (a square with one diagonal). Write down
$\mat{A}$, $\mat{D}$ and $\mat{L}$, verify $\mat{L}\vect{1} = \vect{0}$, and
compute $\vect{x}\T\mat{L}\vect{x}$ for $\vect{x} = (1, 0, 0, 1)\T$ both from
the matrix and from the sum over the edges. Find the four eigenvalues (guess
eigenvectors orthogonal to $\vect{1}$ and check the trace), and give
$\lambda_2$, the time constant of consensus, and the two step-size bounds
$1/d_{\max}$ and $2/\lambda_4$.
\end{exercise}

\begin{exercise}[\difficulty{1}]
\label{exr:ch23-average}
(a) Show that the discrete protocol \cref{eq:ch23-discrete} conserves the
average for every $\eps$, and that the maximum state never grows if
$\eps \le 1/d_{\max}$. (b) Four drones sit on a line at $x = 0$, $1$, $10$
and $11$~m, so that $R_{\mathrm{comm}} = 2$~m splits them into the two
components that never agree with each other. How large must
$R_{\mathrm{comm}}$ be before they agree globally? Give $\lambda_2$ and the
time constant $1/\lambda_2$ at that radius and at $R_{\mathrm{comm}} = 10$~m,
and compare them with the $\lambda_2 = 1$ of \cref{ex:ch23-square}. (c) Prove the identity $\sum_{i<j}\norm{\vect{y}_i - \vect{y}_j}^2 =
n\sum_i\norm{\vect{y}_i - \bar{\vect{y}}}^2$ and deduce
$e_F^2 = \frac{2n}{n-1}e_c^2$ for the complete formation graph; check it on
$e_F(0) = 2.5$~m and $e_c(0) = 1.794$~m of \cref{ex:ch23-square}, remembering
that the example measures $e_F$ on four edges, not six.
\end{exercise}

\begin{exercise}[\difficulty{2}]
\label{exr:ch23-proof}
Prove \cref{thm:ch23-discrete} in full. Show that $\mat{P} = \mat{I} -
\eps\mat{L}$ has the eigenvalues $1 - \eps\lambda_k$ with the eigenvectors of
$\mat{L}$, that the average is conserved, that the disagreement vector stays
orthogonal to $\vect{1}$ and contracts by the factor $\rho$ of
\cref{eq:ch23-discreterate} at every step, and that $\rho < 1$ when the graph
is connected and $\eps < 2/\lambda_n$. Then prove $\lambda_n \le 2d_{\max}$
with Gershgorin's theorem and conclude that $\eps < 1/d_{\max}$ suffices.
Where exactly does connectedness enter?
\end{exercise}

\begin{exercise}[\difficulty{2}]
\label{exr:ch23-directed}
Three drones communicate over one-way links: drone~2 hears drone~1 and
drone~3 hears drone~2, nobody hears anybody else. (a) Write $\mat{L}$ for the
convention that $a_{ij} = 1$ when $i$ hears $j$, check that it has a directed
spanning tree, solve $\dot{\vect{x}} = -\mat{L}\vect{x}$ and show that all
states converge to $x_1(0)$; identify the left eigenvector $\vect{w}$. (b) Add
the link by which drone~1 hears drone~3; show that the digraph is now balanced
and that the states converge to the average. (c) Replace the links by ``drone~2
hears drones~1 and~3'' only; show that there is no spanning tree, compute the
limit of each state, and explain in words why no protocol can reach consensus
on this graph.
\end{exercise}

\begin{exercise}[\difficulty{2}]
\label{exr:ch23-inconsistent}
(a) Sum \cref{eq:ch23-formation} over all drones and show that the centroid
moves with the constant velocity $-\frac{1}{n}\sum_i\sum_{j \in N_i}
\vect{d}_{ij}$; conclude that it is stationary whenever $\vect{d}_{ji} =
-\vect{d}_{ij}$ on every edge, and give an example that is antisymmetric on no
edge at all and still leaves the centroid fixed. (b) With antisymmetric displacements, show that the
controller is gradient descent on $J(\pos) = \frac12\sum_{\set{i,j}\in E}
\norm{\pos_j - \pos_i - \vect{d}_{ij}}^2$ and that $J$ decreases monotonically.
For a triangle with $\vect{d}_{12} = \vect{d}_{23} = (1, 0)\T$ and
$\vect{d}_{31} = (-1, 0)\T$ (the vectors do not close), find the configuration
the swarm settles in and its residual formation error. (c) Explain why
computing $\vect{d}_{ij}$ from offsets makes both problems impossible.
\end{exercise}

\begin{exercise}[\difficulty{2}]
\label{exr:ch23-range}
(a) Show that the set of pairs $(\pos_i, \pos_j)$ with $\norm{\pos_i - \pos_j}
\le R_{\mathrm{comm}}$ is convex and write the constraint for one kept edge and
one step of the MPC horizon of \cref{ch:ch21}. (b) Show that the $Q$ linear
inequalities $\vect{n}_q\T(\pos_i - \pos_j) \le R_{\mathrm{comm}}\cos(\pi/Q)$,
with $\vect{n}_q$ equally spaced unit vectors, describe a polygon inside the
disc, so that a plan that satisfies them keeps the edge in range; how much
range is given away for $Q = 8$? (c) Compute the derivative of the barrier
\cref{eq:ch23-barrier} and check that the resulting velocity is continuous at
$\ell_{\mathrm{act}}$. (d) Explain why $\lambda_2$ alone cannot tell which
edge is about to break, and what to log instead.
\end{exercise}

\begin{exercise}[\difficulty{3}]
\label{exr:ch23-secondorder}
(a) Derive the modal equation \cref{eq:ch23-modes2} from
\cref{eq:ch23-second} without pinning, prove the stability condition of
\cref{prop:ch23-second}, give the oscillation frequencies for the graph of
\cref{ex:ch23-square} with $k_p = 1$ and $k_v = k_d = 0$, and find the $k_v$
that makes the slowest mode critically damped for $k_d = 0$. (b) For the
first-order leader--follower controller without feed-forward and a reference
moving at the constant velocity $v\vect{e}$, derive the steady-state error
$-v\,\bigl((k_c\mat{L} + k_l\mat{B})\inv\vect{1}\bigr)\otimes\vect{e}$ of
\cref{prop:ch23-leader}. Evaluate it for the graph of \cref{ex:ch23-square}
with drone~1 pinned, $k_c = k_l = 1$ and $v = 1$~m/s: how far does each drone
trail its slot, and what formation error remains?
\end{exercise}

\begin{exercise}[\difficulty{3} Coding]
\label{exr:ch23-coding}
This is the Week~11 coding exercise of the study plan (\cref{ch:appA}):
simulate formation control with a communication radius and measure the
formation error during an avoidance manoeuvre. Work in the API of
\code{code/ch23\_consensus.py} but write the functions yourself.
(a) Implement the radius graph, the Laplacian and $\lambda_2$, and reproduce
the spectrum $\set{0, 1, 3, 4}$ of \cref{ex:ch23-square}. (b) Implement
\cref{alg:ch23-formation} for the whole swarm and reproduce
\cref{tab:ch23-formation} and the final square; then pin drone~1 to the
waypoint path of \cref{fig:ch23-leader} with and without feed-forward and
compare the formation errors after $20$~s. (c) Send the square of
\cref{fig:ch23-avoidance} along a straight reference at $1$~m/s with
$R_{\mathrm{comm}} = 3.2$~m and let an intruder cross its path; add the safety
repulsion of \cref{ch:ch15} (or \orca from \cref{ch:ch13} if you have it) and
the range barrier \cref{eq:ch23-barrier}; log at every period the formation
error, the shortest and longest edge, the distance to the intruder,
$\lambda_2$ and the number of communication violations, and plot them.
Reproduce the peak error of $0.36$~m, the recovery below $0.05$~m about $3$~s
later, and a longest edge below $R_{\mathrm{comm}}$. (d) Reduce
$R_{\mathrm{comm}}$ to $2.5$~m so that the diagonals of the square
($2.83$~m) are out of range and the graph switches between the $4$-cycle and
the complete graph during the manoeuvre; report what happens to $\lambda_2$,
to the recovery time and to the violations, and then remove the range barrier
and show a run in which the formation splits. (e) Optional: replace the
formation term by the flocking controller \cref{eq:ch23-flock} and check that
the velocities align and that the group stays together only with the
navigation term switched on.
\end{exercise}
